\section{Observer-Declared Timeline Model}
\label{sec:model}

An observer-declared timeline is a deterministic projection over accepted
facts from a declared perspective. The model is intentionally small; it is a
surface contract, not a full database schema.

\subsection{Facts and sources}

A \emph{fact} is an event, record, manifest, receipt, payload hash, causal
edge, or verification result that the system treats as evidence. A fact has a
source, provenance, and a local identity. A fact source may be a local runtime,
remote runtime, repository, release artifact, adapter, imported bundle, or
external system.

The view does not need to accept every fact the product has ever seen. It must
declare which source ranges, heads, cursors, manifests, watermarks, or
freshness boundaries were accepted for the current projection.

\subsection{Observer}

An \emph{observer} is the perspective whose view is being shown. It may be a
human, agent, runtime location, service, or product view. The observer is not
necessarily the author of the facts. It is the subject from whose accepted
fact set and projection policy the timeline is derived.

\subsection{Projection policy}

A projection policy orders accepted facts for display, replay, export, or
agent consumption. Causal constraints dominate projection policy: if fact
$a$ is known to causally precede fact $b$, a valid projection must not place
$b$ before $a$. Concurrent or causally unrelated facts may be ordered by an
observer policy, source priority, local sequence, or deterministic
tie-breaker, but the policy must be inspectable.

\subsection{Degraded evidence}

A timeline may be useful even when evidence is incomplete. Instead of hiding
that incompleteness, the view should declare degraded states such as missing
causality, incomplete accepted ranges, stale remote facts, redacted payloads,
schema mismatch, unverifiable source, or unresolved conflict. A degraded view
is preferable to an overconfident total order.

\subsection{Minimal surface}

The minimal observer-declared surface contains:

\begin{itemize}[leftmargin=*]
  \item a stable view identifier;
  \item the view subject, such as timeline, history, replay, sync, ordering,
        or mixed-source work state;
  \item observer identity and kind;
  \item accepted fact sources and ranges;
  \item projection policy version, causal dominance, and tie-breaker;
  \item causal constraints;
  \item degraded evidence states; and
  \item verification status.
\end{itemize}

This is close to the KFD-4 observer-perspective schema published by the KFD
package. The schema is a vocabulary and verification boundary, not a claim that
an adopter's runtime implementation is correct merely because it imports the
schema.
